\begin{kbvrule}
\noindent{\kbvsecfont ಸೂತ್ರ :}\hspace{0.5em}ಸಂಧಿಕಾರ್ಯ ನಡೆಯುವಾಗ ಅ, ಆ ಕಾರಗಳಿಗೆ ಕರಳು ಪರವಾದರೆ “ಏ” ಕಾರವೂ, ಉ, ಊ ಕಾರಗಳು ಪರವಾದ ಕಾರವು ಋಕಾರವು ಪರವಾದರೆ “ಅರ್” ಕಾರವೂ ಆದೇಶಗಳಾಗಿ ಬರುವುವು. ಅದು ಗುಣಸಂಧಿ ಎನಿಸುವುದು.
\end{kbvrule}

\vspace{4pt plus 4pt minus 1pt}\noindent{\kbvsecfont\fontsize{11pt}{14.5pt}\selectfont\color{kbvaccentdark} (೩)\hspace{0.45em}ವೃದ್ಧಿ ಸಂಧಿ :}\hspace{0.5em}“ಏಕೈಕ” — ಈ ಶಬ್ದವನ್ನು ಬಿಡಿಸಿ ನೋಡೋಣ.\par

\begin{kbvderiv}
\vspace{4pt plus 3pt minus 1pt}
\begin{center}\begin{adjustbox}{max width=\textwidth}
{\renewcommand{\arraystretch}{1.22}\setlength{\tabcolsep}{4pt}
\begin{tabular}{l c l}
ಏಕೈಕ & {\color{kbvgrey}$=$} & ಏಕ + ಏಕ \\
 & {\color{kbvgrey}$=$} & ಏಕ್ + (ಅ + ಏ) + ಕ \\
\end{tabular}}
\end{adjustbox}\end{center}
\vspace{4pt plus 3pt minus 1pt}
\end{kbvderiv}

\indent ಇಲ್ಲಿ ಅ ಕಾರಕ್ಕೆ ಏ ಕಾರ ಪರವಾಗಿದೆ. ಇಲ್ಲಿ ಅ, ಏ ಗಳ ಸ್ಥಾನದಲ್ಲಿ “ಐ” ಕಾರ ಆದೇಶವಾಗಿ ಬಂದಿದೆ.\par

\indent  = ಏಕೈಕ (“ಕ್” — ಐಕ್)\par

\indent ತೀರ್ಮಾನ : ಅಕಾರಕ್ಕೆ ಏಕಾರ ಪರವಾದಾಗ ಐಕಾರ ಆದೇಶವಾಗಿ ಬರುವುದು.\par

\indent ಇದೇ ರೀತಿಯಾಗಿ,\par

\begin{kbvderiv}
\vspace{4pt plus 3pt minus 1pt}
\begin{center}\begin{adjustbox}{max width=\textwidth}
{\renewcommand{\arraystretch}{1.22}\setlength{\tabcolsep}{4pt}
\begin{tabular}{l c l}
ಅಷ್ಟ + ಐಶ್ವರ್ಯ & {\color{kbvgrey}$=$} & ಅಷ್ಟೈಶ್ವರ್ಯ (ಅ + ಐ = ಐ) \\
ವನ + ಓಷಧಿ & {\color{kbvgrey}$=$} & ವನೌಷಧಿ (ಅ + ಓ = ಔ) \\
ಮಹಾ + ಔದಾರ್ಯ & {\color{kbvgrey}$=$} & ಮಹೌದಾರ್ಯ (ಆ + ಔ = ಔ) \\
ಭಾವ + ಐಕ್ಯತೆ & {\color{kbvgrey}$=$} & ಭಾವೈಕ್ಯತೆ (ಅ + ಐ = ಐ) \\
\end{tabular}}
\end{adjustbox}\end{center}
\vspace{4pt plus 3pt minus 1pt}
\end{kbvderiv}

\kbvsechead{ಈ ಉದಾಹರಣೆಗಳಲ್ಲಿ ಪೂರ್ವದ ಸ್ವರವು ಅ ಅಥವಾ ಆ ಆಗಿರುವುದು :}

\indent ಅವುಗಳಿಗೆ ಪರವಾಗಿರುವ ಸ್ವರಗಳು ಏ, ಐ, ಓ, ಔ ಆಗಿರುವುವು. ಅ, ಆ ಗಳಿಗೆ ಏ, ಐ ಗಳು ಪರವಾದಾಗ “ಐ” ಕಾರವೂ ಓ, ಔ ಗಳು ಪರವಾದಾಗ “ಔ” ಕಾರವೂ ಆದೇಶಗಳಾಗಿ ಬಂದಿರುವುವು. ಈ ತೆರನಾದ ಸಂಧಿಯು ವೃದ್ಧಿ ಸಂಧಿಯೆನಿಸುವುದು.\par
